%!TeX program = pdflatex

\documentclass[11pt,letterpaper]{report}

\usepackage[T1]{fontenc}
\usepackage[strict,autostyle]{csquotes}
\usepackage[USenglish]{babel}
\usepackage{microtype}
\usepackage{datetime}
\usepackage{tabto}
\usepackage{setspace}
\usepackage{hyperref}
\usepackage{geometry}
\usepackage{enumitem}
\usepackage{titlesec}
\usepackage{xcolor}

\newcommand{\myname}{Zaibei (Eric) Li}

\usepackage{ebgaramond}
\usepackage{tgheros}

\newcommand{\listtabwidth}{1.9cm}
\newcommand{\namefont}[1]{{\normalfont\bfseries\Huge{#1}}}

\SetTracking{encoding=*, family=\sfdefault}{30}
\titleformat{\section}{\lsstyle\sffamily\small\bfseries\uppercase}{}{}{}{}
\titlespacing{\section}{0pt}{22pt plus 4pt minus 4pt}{6pt plus 2pt minus 2pt}

\titleformat{\subsection}{\lsstyle\sffamily\footnotesize\bfseries}{}{}{}{}
\titlespacing{\subsection}{0pt}{14pt plus 4pt minus 4pt}{4pt plus 2pt minus 2pt}

\geometry{body={6.7in, 9.2in},
    left=0.9in,
    top=0.8in}

\setlength\parindent{0em}
\setstretch{0.94}

\newcommand{\listitemspace}{0.18em}

\renewenvironment{itemize}
{\begin{list}{}{\setlength{\leftmargin}{0em}
                \setlength{\parskip}{0em}
                \setlength{\itemsep}{\listitemspace}
                \setlength{\parsep}{\listitemspace}}}
{\end{list}}

\TabPositions{\listtabwidth}
\newlist{tablist}{description}{3}
\setlist[tablist]{leftmargin=\listtabwidth,
    labelindent=0em,
    topsep=0em,
    partopsep=0em,
    itemsep=\listitemspace,
    parsep=\listitemspace,
    font=\normalfont}

\newdateformat{monthyeardate}{\monthname[\THEMONTH] \THEYEAR}

\hypersetup{
    colorlinks  = true,
    urlcolor    = black,
    citecolor   = black,
    linkcolor   = black,
    pdfauthor   = {\myname},
    pdfkeywords = {software engineering, research engineering, multimodal learning analytics, IoT, HCI},
    pdftitle    = {\myname: Industry Resume},
    pdfsubject  = {Resume},
    pdfpagemode = UseNone
}

\begin{document}
\raggedright{}

\namefont{\myname}

\vspace{0.8em}
\begin{minipage}[t]{0.62\textwidth}
    Department of Computer Science and Science Education\\
    University of Copenhagen
\end{minipage}
\begin{minipage}[t]{0.37\textwidth}
    \flushright{}
    \href{mailto:zali@di.ku.dk}{zali@di.ku.dk} \\
    +45 91858214 \\
    \href{https://lizaibeim.github.io}{lizaibeim.github.io}
\end{minipage}

\section*{Summary}

Ph.D. researcher in Computer Science building data-driven tools for multimodal learning analytics, human-AI collaboration, and interactive systems. Experience shipping open-source research software, prototyping IoT and audio-based sensing workflows, and turning research ideas into usable tools, dashboards, and empirical studies.

\section*{Technical Skills}

Python, C\#, Unity, LaTeX; multimodal data collection, IoT prototyping, audio analysis, dashboard design, data visualization, human-computer interaction, human-AI collaboration, research prototyping, academic writing.

\section*{Experience}

\begin{tablist}
    \item[2024--] \tab{}\textbf{Doctoral Researcher}, University of Copenhagen, Copenhagen, Denmark
\end{tablist}

\begin{itemize}
    \item Developed open-source tooling for multimodal learning analytics, including \textbf{OpenMMLA}, an IoT-based data collection toolkit for classroom and collaboration research.
    \item Designed and implemented research prototypes spanning sensing, audio analysis, and interactive dashboards to support data collection and post-hoc analysis in authentic learning settings.
    \item Published first-author and co-authored work at LAK, EC-TEL, ICALT, and related venues; received the \textbf{Best Short Paper Award} at LAK 2025.
\end{itemize}

\begin{tablist}
    \item[2023--2024] \tab{}\textbf{Research Assistant}, University of Copenhagen, Copenhagen, Denmark
\end{tablist}

\begin{itemize}
    \item Supported research and system prototyping for multimodal learning analytics and smart learning environments.
    \item Contributed to voice and audio sensing workflows, data collection infrastructure, and early-stage platform design for mBox-related projects.
\end{itemize}

\begin{tablist}
    \item[2024--] \tab{}\textbf{Teaching Assistant}, University of Copenhagen
\end{tablist}

\begin{itemize}
    \item Assisted with \textit{Introduction to Python Programming and Data Science}, supporting undergraduate teaching and hands-on programming exercises.
\end{itemize}

\section*{Selected Projects}

\begin{tablist}
    \item[OpenMMLA] \tab{}Open-source multimodal data collection toolkit for learning analytics. Focused on IoT-based sensing and research workflow support. \href{https://github.com/ucph-ccs/OpenMMLA}{GitHub}
    \item[MotionMatching] \tab{}Real-time motion matching system in Unity and C\# for interactive character motion generation. \href{https://github.com/lizaibeim/motion-matching}{GitHub}
    \item[CasperFFG] \tab{}Python simulation of CasperFFG consensus combined with Proof-of-Work for blockchain experimentation. \href{https://github.com/lizaibeim/casper-ffg}{GitHub}
\end{tablist}

\section*{Selected Publications}

\begin{tablist}
    \item[2025] \tab{}\textbf{Z. Li}, S. Yamaguchi, and D. Spikol. \enquote{OpenMMLA: an IoT-based Multimodal Data Collection Toolkit for Learning Analytics.} LAK '25.
    \item[2024] \tab{}\textbf{Z. Li}, M. T. Jensen, A. Nolte, and D. Spikol. \enquote{Field report for Platform mBox: Designing an Open MMLA Platform.} LAK '24.
    \item[2024] \tab{}D. Spikol, \textbf{Z. Li}, A. Nolte, A. Ohsaki, and K. Rapur. \enquote{Investigating Hackathons with Collaboration Analytics.} ICGJ '24.
\end{tablist}

\section*{Education}

\begin{tablist}
    \item[Ph.D.] \tab{}Computer Science, University of Copenhagen, 2027 (expected)
    \item[M.S.] \tab{}Computer Science, University of Copenhagen, 2022
    \item[B.S.] \tab{}Information Technology, Hong Kong Polytechnic University, 2019
    \item[Exchange] \tab{}Computer Science, KAIST, 2017
\end{tablist}

\section*{Awards}

\begin{tablist}
    \item[2025] \tab{}Best Short Paper Award, LAK '25
    \item[2024--2027] \tab{}Ph.D. Fellowship, University of Copenhagen
\end{tablist}

\begin{center}
    \vfill
    Updated \monthyeardate\today
\end{center}

\end{document}
